\subsection{Example: Jelly-fish busters}
\label{sec:jellyfish-busters}
\slidesrc{BDI architecture: 9--18}

A one-dimensional world of four cells \code{a}, \code{b}, \code{c}, \code{d}. The agent
(SpongeBob in the slides) starts in \code{a}, a jellyfish is in \code{b}, the bin is in
\code{d}. Goal: catch the jellyfish and throw it into the bin.

\begin{center}
\begin{tikzpicture}[font=\small]
  \foreach \x/\c in {0/a,2.0/b,4.0/c,6.0/d} \draw[thick, fill=nbblue!8] (\x,0) rectangle ++(1.9,1.0) node[pos=0.15, font=\scriptsize] {\c};
  \node at (0.95,0.5) {\textbf{agent}};
  \node at (2.95,0.5) {\textcolor{nbpurple}{\textbf{jellyfish}}};
  \node at (6.95,0.5) {\textcolor{nbgreen!60!black}{\textbf{bin}}};
\end{tikzpicture}
\end{center}

\inputminted{prolog}{sections/bdi-agentspeak/jellyfish-busters/assets/jellyfish.asl}

\begin{itemize}[nosep]
  \item \textbf{Initial beliefs} (slide 9): the map (\code{adjacent/2}) and the positions
        (\code{location/2}).
  \item \textbf{Plan 1} (slide 10): \emph{if} I am where the jellyfish is, \emph{then} pick
        it, move to the bin, drop it. \code{!pick} and \code{!drop} are subgoals.
  \item \textbf{Plan 2} (slide 11): \emph{if} the jellyfish is not here, \emph{then} move to
        an adjacent cell and try again: a \emph{recursive} goal, like a recursive Prolog
        predicate.
  \item \textbf{Initial goal} (slide 12): \code{!bust\_jellyfishes}, which generates the
        external event \code{+!bust\_jellyfishes} at start-up.
\end{itemize}

\begin{nbexample}{the run, step by step (slides 13--17)}
\small
\begin{tabularx}{\linewidth}{@{}c >{\raggedright\arraybackslash}p{2.9cm} X >{\raggedright\arraybackslash}p{2.7cm}@{}}
\toprule
\textbf{Step} & \textbf{Event} & \textbf{Plan selection} & \textbf{Effect} \\
\midrule
1 & \code{+!bust_jellyfishes} (external) &
  Relevant: plans 1 and 2. Plan 1: context \code{location(myself,X)} gives \code{X = a},
  but \code{location(jellyfish, a)} is false: not applicable. Plan 2: \code{X = a},
  \code{not location(jellyfish, a)} true, \code{adjacent(a, Z)} gives \code{Z = b}:
  applicable. & new intention; executes \code{move_to(b)} \\
2 & --- & the body continues & a percept updates the agent's position to \code{b} \\
3 & \code{+!bust_jellyfishes} (internal) &
  posted by plan 2's body. Plan 1: \code{X = b}, \code{location(jellyfish, b)} true,
  \code{location(bin, Y)} gives \code{Y = d}: applicable. Plan 2: not applicable
  (\code{not location(jellyfish, b)} is false). & plan 1 pushed on the same intention \\
4 & \code{+!pick}\newline\code{(jellyfish,b)} & handled by a \code{pick} plan (not shown) & jellyfish
  picked up (slide 15) \\
5 & --- & \code{move_to(d)} & agent carries it to \code{d} (slide 16) \\
6 & \code{+!drop}\newline\code{(jellyfish,d)} & handled by a \code{drop} plan (not shown) & jellyfish in
  the bin (slide 17); plans popped: intention completed \\
\bottomrule
\end{tabularx}
\end{nbexample}

\begin{center}
\begin{tikzpicture}[font=\scriptsize\ttfamily, >=Stealth,
  f/.style={draw, thick, minimum width=3.6cm, minimum height=0.5cm, fill=#1}]
  \node[font=\small\rmfamily] at (0,1.0) {after step 1};
  \node[f=nbblue!10] at (0,0.3) {plan 2 [X=a, Z=b]};
  \node[font=\small\rmfamily] at (4.6,1.6) {after step 3};
  \node[f=nborange!12] at (4.6,0.9) {plan 1 [X=b, Y=d]};
  \node[f=nbblue!10] at (4.6,0.3) {plan 2 [X=a, Z=b]};
  \node[font=\small\rmfamily] at (9.2,2.2) {during step 4};
  \node[f=nbpurple!10] at (9.2,1.5) {pick plan};
  \node[f=nborange!12] at (9.2,0.9) {plan 1 [X=b, Y=d]};
  \node[f=nbblue!10] at (9.2,0.3) {plan 2 [X=a, Z=b]};
  \draw[thick] (-1.9,0.02) -- (1.9,0.02); \draw[thick] (2.7,0.02) -- (6.5,0.02); \draw[thick] (7.3,0.02) -- (11.1,0.02);
\end{tikzpicture}
\captionof{figure}{The single intention of the run as a stack of partially instantiated plans:
each subgoal pushes a new plan on top; when the top plan finishes it is popped and the plan
below continues.}
\end{center}

\begin{nbwarning}[What the slides leave implicit]
\begin{itemize}[nosep]
  \item The program has no plans for \code{+!pick(\ldots)} and \code{+!drop(\ldots)}: in a
        complete program they would be plans (or be written as actions \code{pick(\ldots)},
        \code{drop(\ldots)} without \code{!}). An achievement goal with no relevant plan
        \emph{fails}, and so does the intention.
  \item The belief \code{location(myself, \ldots)} must be updated after each move: in
        Jason this is done by the \emph{environment} through percepts (or by the plan with
        \code{-+location(myself, Z)}). Otherwise plan 2 would move to \code{b} forever.
  \item If the jellyfish were in no reachable cell, plan 2 would reach \code{d}, find no
        \code{adjacent(d, \_)}, and no plan would be applicable: goal failure (a
        \code{-!bust\_jellyfishes} plan could handle it).
\end{itemize}
\end{nbwarning}

\begin{nbwhy}[Jason (slide 18)]
\textbf{Jason} (Bordini, H\"ubner, Wooldridge, \emph{Programming Multi-Agent Systems in
AgentSpeak Using Jason}, Wiley 2007) is the best-known interpreter of an extended
AgentSpeak, written in Java. Slide 18 shows its ``Mars world'' example: robots \code{r1} and
\code{r2} on a grid, collecting garbage \code{G} and carrying it to an incinerator, with the
mind inspector showing beliefs, events and intentions at each reasoning cycle. Jason adds to
AgentSpeak(L) communication between agents (via speech-act based messages,
Section~\ref{sec:agent-communication}), annotations on beliefs (e.g.\ \code{[source(percept)]}),
environments in Java, and much more; it is the platform of the course labs.
\end{nbwhy}
